\subsection{Extending strictly positive potentials}
\label{sec:strict-positive-potential-extension}
We keep the compact K\"ahler manifold $(X,\omega)$, but now we extend potentials instead of sections. Let $Y\subset X$ be a closed complex submanifold, and let $\varphi\in\PSH(Y,\omega|_Y)$ be such that $\omega|_Y+\ddc\varphi$ is a K\"ahler current on $Y$, that is, dominates a positive multiple of $\omega|_Y$. We look for $\widetilde\varphi\in\PSH(X,\omega)$ that restricts exactly to $\varphi$ and for which $\omega+\ddc\widetilde\varphi$ is again a K\"ahler current (\cref{thm:holder-positive-potential-extension}). Near $Y$, local extensions are obtained by pulling $\varphi$ back by a holomorphic projection onto $Y$ and adding a large multiple of a function vanishing to second order along $Y$. The difficulty is to glue them, since two projections give different pullbacks. We glue by a maximum after lowering each local extension by a cutoff penalty near the edge of its chart. The maximum is then $\omega$-plurisubharmonic near $Y$ and keeps the prescribed values on $Y$, provided that near every point of $Y$ on the edge of a chart the penalized extension from that chart is dominated by the extension from another chart. This comparison is obtained by adding to $\mathrm e^\varphi$ the exponential of a small logarithmic pole along $Y$: when $\mathrm e^\varphi$ is H\"older continuous, the difference between two pullbacks of $\mathrm e^\varphi$ is small compared with this exponential near $Y$. By \cref{lem:background-logarithmic-sum} below, the logarithm of such a sum is again quasi-plurisubharmonic.

H\"older continuity on a manifold is defined in charts. Let $M$ be a smooth manifold of dimension $d$. A function $f\colon M\to\mathbb R$ is \emph{H\"older continuous} if every point of $M$ has a coordinate neighborhood $V$, with coordinates $x\colon V\to\mathbb R^d$, on which there are constants $C\geq0$ and $\alpha\in(0,1]$ such that
\[
 |f(a)-f(b)|\leq C|x(a)-x(b)|^\alpha\qquad\text{for all }a,b\in V.
\]
We then say that $f$ is H\"older continuous with \emph{exponent} $\alpha$ on $V$, in the coordinates $x$. Such a function is continuous on $V$, and hence bounded near each point of $V$.
\index{H\"older continuity!on manifolds}

\begin{lemma}\label{lem:manifold-holder-operations}
Let $M$ and $M'$ be smooth manifolds.
\begin{enumerate}
\item Let $f$ be H\"older continuous with exponent $\alpha$ on a coordinate neighborhood $V$ in the coordinates $x$, let $p\in V$, and let $y$ be coordinates on a neighborhood of $p$. Then $f$ is H\"older continuous with exponent $\alpha'$ in the coordinates $y$ on some neighborhood of $p$, for every $\alpha'\in(0,\alpha]$.
\item If $M$ is compact and $f$ is H\"older continuous on $M$, then there are $\alpha\in(0,1]$ and finitely many coordinate neighborhoods covering $M$ on each of which $f$ is H\"older continuous with exponent $\alpha$.
\item Smooth functions are H\"older continuous. If $f$ and $g$ are H\"older continuous on $M$, then so are $f+g$, $fg$ and $\max\{f,g\}$. If $f\geq0$ is H\"older continuous and $0<\gamma\leq1$, then $f^\gamma$ is H\"older continuous.
\item If $\Phi\colon M'\to M$ is a smooth map and $f$ is H\"older continuous on $M$, then $f\circ\Phi$ is H\"older continuous on $M'$.
\end{enumerate}
\end{lemma}
\begin{proof}
(1) Choose a closed ball $N$ in the $y$-coordinates, centered at $y(p)$, whose preimage lies in $V$. The map $x\circ y^{-1}$ is smooth on the convex compact set $N$, so by the mean value inequality it is Lipschitz on $N$ with some constant $L$. Moreover $N$ has some diameter $D$, and $t^\alpha\leq D^{\alpha-\alpha'}t^{\alpha'}$ for $0\leq t\leq D$. Hence, for $a,b$ in the preimage of $N$, which is a neighborhood of $p$,
\[
 |f(a)-f(b)|\leq C|x(a)-x(b)|^\alpha\leq CL^\alpha|y(a)-y(b)|^\alpha\leq CL^\alpha D^{\alpha-\alpha'}|y(a)-y(b)|^{\alpha'}.
\]

(2) By (1), applied with $y=x$, every point of $M$ has a coordinate neighborhood on which $f$ is H\"older continuous with some exponent, and on which it remains so with every smaller exponent. Finitely many of these neighborhoods cover $M$, and the smallest of their exponents serves on all of them.

(3) A smooth function is Lipschitz on every closed coordinate ball, by the mean value inequality, so it is H\"older continuous with exponent one. Let $f$ and $g$ be H\"older continuous and $p\in M$. By (1) there are coordinates $x$ on a neighborhood $V$ of $p$ and $\alpha\in(0,1]$ such that $f$ and $g$ are H\"older continuous with exponent $\alpha$ on $V$ in these coordinates, say with the constant $C$; shrinking $V$, we may assume that $|f|,|g|\leq B$ on $V$. For $a,b\in V$ we have $|(f+g)(a)-(f+g)(b)|\leq2C|x(a)-x(b)|^\alpha$,
\[
 |f(a)g(a)-f(b)g(b)|\leq|f(a)|\,|g(a)-g(b)|+|g(b)|\,|f(a)-f(b)|\leq2BC|x(a)-x(b)|^\alpha,
\]
and $|\max\{f,g\}(a)-\max\{f,g\}(b)|\leq\max\{|f(a)-f(b)|,|g(a)-g(b)|\}$. For the power, note first that $(s+t)^\gamma\leq s^\gamma+t^\gamma$ for $s,t\geq0$: if $s+t>0$, divide by $(s+t)^\gamma$ and use $\lambda^\gamma\geq\lambda$ for $\lambda\in[0,1]$. Hence, if $s\geq t\geq0$, then $s^\gamma\leq(s-t)^\gamma+t^\gamma$, that is, $|s^\gamma-t^\gamma|\leq|s-t|^\gamma$ for all $s,t\geq0$. If $f\geq0$ has exponent $\alpha$ on $V$ with the constant $C$, it follows that $|f(a)^\gamma-f(b)^\gamma|\leq C^\gamma|x(a)-x(b)|^{\alpha\gamma}$.

(4) Let $p'\in M'$, let $V$ be a coordinate neighborhood of $\Phi(p')$, with coordinates $x$, on which $f$ is H\"older continuous with exponent $\alpha$ and constant $C$, and let $N$ be a closed coordinate ball around $p'$, in coordinates $x'$, with $\Phi(N)\subset V$. The map $x\circ\Phi$, written in the coordinates $x'$, is smooth on the convex compact set $N$, hence Lipschitz on $N$ with some constant $L$. Therefore $|f(\Phi(a))-f(\Phi(b))|\leq CL^\alpha|x'(a)-x'(b)|^\alpha$ for $a,b\in N$.
\end{proof}

\begin{lemma}\label{lem:background-logarithmic-sum}
Let $U$ be an open subset of a complex manifold, let $\theta$ be a smooth closed real $(1,1)$-form on $U$, and let $a,b\in\PSH(U,\theta)$. Then
\[
 \log(\mathrm e^a+\mathrm e^b)\in\PSH(U,\theta),
\]
where $\mathrm e^{-\infty}=0$ and $\log0=-\infty$, so that this function equals $-\infty$ precisely at the common poles of $a$ and $b$. Moreover
\[
 \max\{a,b\}\leq\log(\mathrm e^a+\mathrm e^b)\leq\max\{a,b\}+\log2.
\]
\end{lemma}
\begin{proof}
Both assertions are local. Let $B\Subset U$ be a coordinate ball, let $s$ be a smooth real function on $B$ with $\ddc s=\theta$ (\cref{lem:smooth-potential-on-ball}), and let $B'\Subset B$ be a smaller concentric ball. Then $a+s$ and $b+s$ are plurisubharmonic on $B$, by the remarks following \cref{def:quasi-psh-and-theta-psh}, and
\[
 \log(\mathrm e^a+\mathrm e^b)+s=H(a+s,b+s),\qquad H(x,y)=\log(\mathrm e^x+\mathrm e^y).
\]
The function $H$ is smooth on $\mathbb R^2$ and increasing in each variable: its partial derivatives are $p$ and $1-p$, where $p=\mathrm e^x/(\mathrm e^x+\mathrm e^y)$, and its Hessian matrix is $p(1-p)\left(\begin{smallmatrix}1&-1\\-1&1\end{smallmatrix}\right)$, so $H$ is also convex. If $f,g$ are smooth plurisubharmonic functions, the chain rule gives
\[
 \ddc H(f,g)=p\,\ddc f+(1-p)\,\ddc g
      +p(1-p)\,\frac{\mathrm i}{2\pi}\,\partial(f-g)\wedge\bar\partial(f-g)\geq0,
\]
with $p$ evaluated at $(f,g)$; hence $H(f,g)$ is plurisubharmonic, by the criterion for smooth functions stated after \cref{thm:psh-distribution-representative}. In general, let $f_k$ and $g_k$ be the regularizations of $a+s$ and $b+s$ by radial mollifiers whose radii decrease to zero and are smaller than the distance from $B'$ to $\partial B$. They are smooth plurisubharmonic functions on $B'$, and by \eqref{eq:canonical-subharmonic} and \cref{thm:psh-distribution-representative} they decrease to $a+s$ and $b+s$ (see the proof of \cref{prop:bt-positive-recursion}). The function $H$ extends continuously to $[-\infty,\infty)^2$, with values in $[-\infty,\infty)$, and it remains increasing in each variable; the continuity at points with a coordinate $-\infty$ follows from the inequalities proved below. Hence the functions $H(f_k,g_k)$ decrease to $H(a+s,b+s)$. This limit is at least $a+s$, so it is not identically $-\infty$, and it is plurisubharmonic on $B'$ by \cref{thm:psh-decreasing-nets}. Consequently $\log(\mathrm e^a+\mathrm e^b)=H(a+s,b+s)-s$ is quasi-plurisubharmonic on $B'$ and $\theta+\ddc\log(\mathrm e^a+\mathrm e^b)=\ddc H(a+s,b+s)\geq0$ there; since every point of $U$ is the center of such a ball $B'$, the first assertion follows.

For the inequalities fix a point and suppose that $a\geq b$ there; the other case is symmetric. If $a=-\infty$, then $b=-\infty$ as well, and all three expressions equal $-\infty$. Otherwise $\log(\mathrm e^a+\mathrm e^b)=a+\log(1+\mathrm e^{b-a})$, and $0\leq\log(1+\mathrm e^{b-a})\leq\log2$. The same inequalities for $H$, namely $\max\{x,y\}\leq H(x,y)\leq\max\{x,y\}+\log2$, show that $H(x,y)\to-\infty$ when $\max\{x,y\}\to-\infty$, and that $H(x,y)\to y_0$ when $x\to-\infty$ and $y\to y_0\in\mathbb R$, because $H(x,y)=y+\log(1+\mathrm e^{x-y})$; this is the continuity used above.
\end{proof}
\index{quasi-plurisubharmonic function!logarithm of a sum of exponentials}

\begin{theorem}\label{thm:holder-positive-potential-extension}
Let $Y\subset X$ be a closed complex submanifold, and let $\varphi\in\PSH(Y,\omega|_Y)$ satisfy
\[
 \omega|_Y+\ddc\varphi\geq\delta\,\omega|_Y
\]
for some $\delta>0$. Suppose that $\mathrm e^\varphi$, defined to be zero at the poles of $\varphi$, is H\"older continuous on $Y$. Then there are $\widetilde\varphi\in\PSH(X,\omega)$ and $\varepsilon>0$ such that
\[
 \widetilde\varphi|_Y=\varphi,\qquad
 \omega+\ddc\widetilde\varphi\geq\varepsilon\omega.
\]
Moreover $\widetilde\varphi$ is continuous on $X\setminus Y$, and $\mathrm e^{\widetilde\varphi}$ is H\"older continuous on $X$. If $\varphi$ has analytic singularities, then so does $\widetilde\varphi$.
\end{theorem}
\begin{proof}
\emph{Step 0. Reductions.} If $Y=\varnothing$, the function $\widetilde\varphi=0$ and $\varepsilon=1$ have all the stated properties; $0=\log|1|^2$ has analytic singularities. Let $Y\neq\varnothing$. The connected components of $Y$ are open in $Y$, and $Y$ is compact, so there are finitely many of them, say $Y_1,\ldots,Y_N$; each $Y_k$ is closed in $X$, hence a compact connected complex submanifold, and different components may have different dimensions. If some $Y_k$ has dimension $n$, then it is open in $X$, because near each of its points it is a coordinate plane of full dimension; being also closed, it equals $X$, as $X$ is connected. Then $Y=X$, and $\widetilde\varphi=\varphi$ with $\varepsilon=\delta$ has the stated properties. We may therefore assume that every component of $Y$ has positive codimension; components of dimension zero are allowed. Fix $\varepsilon$ with $0<\varepsilon<\min\{\delta/4,1/4\}$. By \cref{lem:manifold-holder-operations}(2), $\mathrm e^\varphi$ is H\"older continuous with one exponent $\alpha\in(0,1]$ on finitely many coordinate neighborhoods covering $Y$; by \cref{lem:manifold-holder-operations}(1), every point $p\in Y$ then has, in every system of coordinates of $Y$ near $p$, a neighborhood on which $\mathrm e^\varphi$ is H\"older continuous with exponent $\alpha$.

\emph{Step 1. A small normal pole.} For each $k$ let $R_k=R_{Y_k}$ be the function of \cref{lem:global-normal-pole} for the submanifold $Y_k$, and let $b_k\geq0$ be the constant called $b$ there. Put
\[
 R=R_1\cdots R_N,\qquad\kappa=b_1+\cdots+b_N.
\]
Then $R\geq0$ is smooth and vanishes exactly on $Y$. The function $\log R=\sum_k\log R_k$ is quasi-plurisubharmonic, satisfies $\ddc\log R\geq-\kappa\omega$, and has analytic singularities by \cref{prop:rational-analytic-operations}, as a sum of functions with analytic singularities. Choose a rational number $c$ with $0<c<\alpha$ and $c\kappa/2\leq1-4\varepsilon$; this is possible because $1-4\varepsilon>0$. Put $F=(c/2)\log R$. Then
\[
 \omega+\ddc F\geq\Bigl(1-\frac{c\kappa}2\Bigr)\omega\geq4\varepsilon\omega.
\]
The function $F$ is smooth on $X\setminus Y$, equals $-\infty$ on $Y$, and has analytic singularities, since $c/2$ is rational. Its exponential $\mathrm e^F=R^{c/2}$ is H\"older continuous on $X$ by \cref{lem:manifold-holder-operations}(3).

We need a lower bound for $R$ near $Y$. Let $p\in Y_k$, and let $(z',w)$ be holomorphic coordinates centered at $p$ in which $Y=\{w=0\}$ near $p$. The function $P=\prod_{l\neq k}R_l$ is smooth and positive near $p$, since the closed sets $Y_l$, $l\neq k$, do not contain $p$. Since $R_k\geq0$ vanishes on $Y_k$, its first derivatives vanish there. Hence the first derivatives of $R=R_kP$ vanish at the points of $Y$ near $p$, and its real Hessian there equals $P$ times the real Hessian of $R_k$. By \cref{lem:global-normal-pole}, the Hessian of $R_k$ at points of $Y_k$ is positive definite in the directions normal to $Y_k$; as a vector $(0,\xi)$ with $\xi\neq0$ is not tangent to $Y=\{w=0\}$, the real Hessian $D_w^2R$ of $w\mapsto R(z',w)$ is positive definite at the points $(z',0)$ near $p$. Since the second derivatives of $R$ are continuous, there are a closed ball $N$ centered at $p$ and a constant $c_0>0$ with $D_w^2R(z',w)[\xi,\xi]\geq2c_0|\xi|^2$ for $(z',w)\in N$ and all $\xi$. If $(z',w)\in N$, the segment $t\mapsto(z',tw)$, $0\leq t\leq1$, lies in $N$, and the function $g(t)=R(z',tw)$ satisfies $g(0)=g'(0)=0$; Taylor's formula with integral remainder gives $R(z',w)=\int_0^1(1-t)g''(t)\,\mathrm dt\geq c_0|w|^2$. Thus
\begin{equation}\label{eq:normal-exponential-lower-bound}
 R\geq c_0|w|^2,\qquad \mathrm e^F=R^{c/2}\geq c_0^{c/2}|w|^c\qquad\text{on }N.
\end{equation}

\emph{Step 2. Local extensions.} Every point of $Y$ has holomorphic coordinates $(z',w)\in\mathbb C^m\times\mathbb C^{n-m}$, centered at that point, in which $Y=\{w=0\}$ on the coordinate domain, where $m<n$ is the dimension of the component of $Y$ through the point; a product of two small balls centered at the origin is then a chart $\widehat W=\widehat B'\times\widehat B''$ with $Y\cap\widehat W=\widehat B'\times\{0\}$. The products $B'\times B''$ of the concentric balls of half the radii cover $Y$, and since $Y$ is compact, finitely many of them suffice. We thus obtain finitely many charts $W_j=B'_j\times B''_j$, relatively compact in charts $\widehat W_j=\widehat B'_j\times\widehat B''_j$ of the same form, with $Y\cap W_j=B'_j\times\{0\}$ and with dimensions $m_j<n$. Let $\pi_j\colon\widehat W_j\to Y$, $\pi_j(z',w)=(z',0)$, be the holomorphic projection. Below we shall shrink the normal factors $B''_j$ several times, keeping them balls centered at the origin; this does not change $Y\cap W_j$, so the charts keep covering $Y$. If a property holds on an open set $O\subset\widehat W_j$ containing the compact set $\overline{B'_j}\times\{0\}$, then this compact set has positive distance from the complement of $O$ in the coordinates, so $O$ contains $W_j$ after the normal factor has been shrunk; we use this observation several times. We claim that there is $A>0$ such that, after such a shrinking,
\[
 a_j=\varphi\circ\pi_j+AR,\qquad b_j=F+AR
\]
belong to $\PSH(W_j,(1-2\varepsilon)\omega)$ for every $j$.

Since $4\varepsilon\leq\delta$, the hypothesis gives $\ddc\varphi\geq-(1-\delta)\omega|_Y\geq-(1-4\varepsilon)\omega|_Y$. Choose a ball $\widetilde B'_j$, concentric with $B'_j$, with $\overline{B'_j}\subset\widetilde B'_j\Subset\widehat B'_j$, and a smooth function $s$ on $\widetilde B'_j\times\{0\}$ with $\ddc s=(1-4\varepsilon)\omega|_Y$ there (\cref{lem:smooth-potential-on-ball}; if $m_j=0$, take $s=0$). Then $\varphi+s$ is plurisubharmonic on $\widetilde B'_j\times\{0\}$, by the remarks following \cref{def:quasi-psh-and-theta-psh}. By \cref{prop:psh-holomorphic-pullback}, $(\varphi+s)\circ\pi_j$ is plurisubharmonic or identically $-\infty$ on the connected set $\widetilde B'_j\times\widehat B''_j$, and the second alternative is excluded because $\pi_j$ maps this set onto $\widetilde B'_j\times\{0\}$, where $\varphi+s$ is not identically $-\infty$. Since $\ddc(s\circ\pi_j)=\pi_j^*\ddc s$ for the smooth function $s$, we obtain that $\varphi\circ\pi_j$ is quasi-plurisubharmonic on $\widetilde B'_j\times\widehat B''_j\supset\overline{W_j}$ and
\[
 \ddc(\varphi\circ\pi_j)\geq-(1-4\varepsilon)\pi_j^*(\omega|_Y).
\]
For $a_j$ it therefore suffices that the smooth form
\[
 Q_j=(1-2\varepsilon)\omega-(1-4\varepsilon)\pi_j^*(\omega|_Y)+A\ddc R
\]
be positive on $W_j$. At a point of $Y$, write the coefficient matrices of real $(1,1)$-forms in blocks with respect to the tangential coordinates $z'$ and the normal coordinates $w$. The form $\pi_j^*(\omega|_Y)$ involves only $\mathrm dz'$ and $\mathrm d\bar z'$, and at a point $(z',0)$ its coefficients are those of $\omega|_Y$ at that point. Hence its mixed and normal blocks vanish, its tangential block at points of $Y$ is the tangential block of $\omega$, and the tangential block of $Q_j$ there is that of $2\varepsilon\omega$, which is positive definite. The function $R$ vanishes on $Y$ together with its first derivatives. Differentiating the identities $R(z',0)=0$ and $\partial R/\partial\bar w_b(z',0)=0$ with respect to $z'$ and $\bar z'$ shows that $\ddc R$ has zero tangential and mixed blocks along $Y$. Its normal block is $(2\pi)^{-1}$ times the matrix $N_R=(\partial^2R/\partial w_a\partial\bar w_b)$, which is positive definite at points of $Y$: for $\xi\in\mathbb C^{n-m_j}\setminus\{0\}$, applying the identity $\Delta=4\partial_t\partial_{\bar t}$ to $t\mapsto R(z',t\xi)$ at $t=0$ gives
\[
 \sum_{a,b}\frac{\partial^2R}{\partial w_a\partial\bar w_b}(z',0)\xi_a\bar\xi_b=\frac14\bigl(D_w^2R(z',0)[\xi,\xi]+D_w^2R(z',0)[\mathrm i\xi,\mathrm i\xi]\bigr)>0
\]
by Step 1. Let $\Omega_M$ and $\Omega_N$ denote the mixed and normal blocks of $(1-2\varepsilon)\omega$, and $\Omega_T$ the tangential block of $\omega$. By the Schur complement criterion for Hermitian block matrices whose upper left block is positive definite, $Q_j$ is positive definite at a point of $Y$ if and only if
\[
 \Omega_N+\frac A{2\pi}N_R-\Omega_M^*\bigl(2\varepsilon\,\Omega_T\bigr)^{-1}\Omega_M
\]
is positive definite there. (If $m_j=0$, there is no tangential block, and the condition reads $\Omega_N+AN_R/(2\pi)>0$.) All blocks depend continuously on the point, $N_R$ is positive definite, and the sets $\overline{B'_j}\times\{0\}$ are compact and finite in number; hence the smallest eigenvalue of $N_R$ on them has a positive lower bound, the other two terms are bounded, and one $A$ achieves this at every point of every $\overline{B'_j}\times\{0\}$. Fix such an $A$. Since the coefficients of $Q_j$ are continuous on $\widehat W_j$ and positive definiteness is an open condition, $Q_j$ is positive definite on an open set containing $\overline{B'_j}\times\{0\}$, and we shrink the normal factor $B''_j$ so that $W_j$ lies in this set. This proves the claim for $a_j$. Along $Y$ the form $\ddc R$ is positive semidefinite, so $A\ddc R+\varepsilon\omega$ is positive definite on an open set containing $\overline{B'_j}\times\{0\}$; after shrinking the normal factors once more we have $A\ddc R\geq-\varepsilon\omega$ on every $W_j$. Together with $\ddc F\geq-(1-4\varepsilon)\omega$ this gives
\[
 (1-2\varepsilon)\omega+\ddc b_j\geq(1-2\varepsilon)\omega-(1-4\varepsilon)\omega-\varepsilon\omega=\varepsilon\omega\geq0,
\]
which proves the claim for $b_j$.

By \cref{lem:background-logarithmic-sum}, applied with $\theta=(1-2\varepsilon)\omega$, the functions
\[
 \varphi_j=\log(\mathrm e^{a_j}+\mathrm e^{b_j})
          =\log(\mathrm e^{\varphi\circ\pi_j}+\mathrm e^F)+AR
\]
belong to $\PSH(W_j,(1-2\varepsilon)\omega)$. On $Y$ we have $F=-\infty$, $R=0$ and $\varphi\circ\pi_j=\varphi$, so $\varphi_j=\varphi$ on $Y\cap W_j$, also at the poles of $\varphi$. Off $Y$ the function $\mathrm e^F$ is positive and continuous and $\mathrm e^{\varphi\circ\pi_j}$ is continuous, so $\varphi_j$ is continuous on $W_j\setminus Y$. Finally $\mathrm e^{\varphi_j}=(\mathrm e^\varphi\circ\pi_j+\mathrm e^F)\,\mathrm e^{AR}$ is H\"older continuous on $W_j$ by \cref{lem:manifold-holder-operations}(3),(4).

\emph{Step 3. Comparison on overlaps.} We claim that for every $p\in Y\cap W_i\cap W_j$ and every $\eta>0$ there is a neighborhood of $p$ on which
\begin{equation}\label{eq:holder-extension-overlap}
 \varphi_i\leq\varphi_j+\eta.
\end{equation}
Work in the coordinates $(z',w)$ of $W_i$, and write $p=(z'_p,0)$. The coordinate $z'$ restricts to coordinates of $Y$ on $Y\cap W_i$, so by Step 0 there are a neighborhood $V_p$ of $p$ in $Y\cap W_i$ and a constant $C_H$ with $|\mathrm e^{\varphi(a)}-\mathrm e^{\varphi(b)}|\leq C_H|z'(a)-z'(b)|^\alpha$ for $a,b\in V_p$. Since $\pi_i$ and $\pi_j$ are continuous and fix $p$, there is a closed ball $N\subset W_i\cap W_j$ centered at $p$ such that $\pi_i(N)\cup\pi_j(N)\subset V_p$ and \eqref{eq:normal-exponential-lower-bound} holds on $N$ in these coordinates. The map $z'\circ\pi_j$ is smooth on the convex compact set $N$, so it is Lipschitz there with some constant $L$. Let $x=(z',w)\in N$, and put $y=(z',0)=\pi_i(x)$. Then $|y-p|\leq|x-p|$, so $y\in N\cap Y$, and $\pi_j(y)=y$ because $\pi_j$ is the identity on $Y\cap W_j$. Hence
\[
 |z'(\pi_i(x))-z'(\pi_j(x))|=|z'(\pi_j(y))-z'(\pi_j(x))|\leq L|y-x|=L|w|,
\]
and therefore
\[
 \bigl|\mathrm e^{\varphi\circ\pi_i}(x)-\mathrm e^{\varphi\circ\pi_j}(x)\bigr|\leq C_HL^\alpha|w|^\alpha.
\]
Since $c<\alpha$, we have $C_HL^\alpha|w|^\alpha\leq(\mathrm e^\eta-1)c_0^{c/2}|w|^c$ whenever $|w|$ is sufficiently small, and then \eqref{eq:normal-exponential-lower-bound} gives $C_HL^\alpha|w|^\alpha\leq(\mathrm e^\eta-1)\mathrm e^{F(x)}$. At such points $x\in N$, including the points of $Y$,
\[
 \mathrm e^{\varphi\circ\pi_i}+\mathrm e^F
 \leq\mathrm e^{\varphi\circ\pi_j}+\mathrm e^\eta\mathrm e^F
 \leq\mathrm e^\eta\bigl(\mathrm e^{\varphi\circ\pi_j}+\mathrm e^F\bigr).
\]
The logarithm is increasing on $[0,\infty)$, with $\log0=-\infty$, so taking logarithms and adding $AR$ to both sides gives \eqref{eq:holder-extension-overlap} on the neighborhood of $p$ formed by the interior points $x$ of $N$ with $|w(x)|$ sufficiently small.

\emph{Step 4. Gluing near $Y$.} Every point of $Y$ has a relatively compact open neighborhood contained in some $W_j$, and finitely many of these neighborhoods cover $Y$; let $W'_j$ be the union of those chosen inside $W_j$. Then $W'_j\Subset W_j$ and the $W'_j$ cover $Y$. Choose open sets $U_j$ and $T_j$ with $W'_j\Subset U_j\Subset T_j\Subset W_j$, for instance sets of points whose coordinate distance from $\overline{W'_j}$ is less than suitable small constants. The compact set $\overline{T_j}\setminus U_j$ is contained in the open set $W_j\setminus\overline{W'_j}$, so by \cref{prop:manifold-cutoffs-stokes} there is a smooth function $f_j\colon X\to[0,1]$ with compact support in $W_j\setminus\overline{W'_j}$ which equals one on a neighborhood of $\overline{T_j}\setminus U_j$; in particular $f_j=0$ on the neighborhood $X\setminus\Supp f_j$ of $\overline{W'_j}$. Fix $\eta>0$ so small that $\eta\ddc f_j\leq\varepsilon\omega$ for every $j$. Since $(1-\varepsilon)\omega+\ddc\varphi_j\geq\varepsilon\omega$, the functions
\[
 \widehat\varphi_j=\varphi_j-\eta f_j
\]
belong to $\PSH(W_j,(1-\varepsilon)\omega)$. On $W=\bigcup_jT_j$ define
\[
 \psi(x)=\max\{\widehat\varphi_j(x):x\in T_j\}.
\]
The set of indices in this maximum changes with $x$. We show that near $Y$ the maximum is nevertheless taken over a fixed set of indices. Let $p\in Y$, and choose $j$ with $p\in W'_j$; then $f_j=0$ near $p$, and $j$ occurs in the maximum near $p$. Let $l$ be an index with $p\in\partial T_l$. Then $p\in\overline{T_l}\setminus U_l$, so $f_l=1$ near $p$. Since $p\in Y\cap W_l\cap W_j$, the inequality \eqref{eq:holder-extension-overlap} with $\eta/2$ in place of $\eta$ gives, near $p$,
\[
 \widehat\varphi_l=\varphi_l-\eta\leq\varphi_j-\eta/2\leq\widehat\varphi_j.
\]
These are inequalities between functions with values in $[-\infty,\infty)$, valid also at common poles. Since there are finitely many indices, there is therefore a neighborhood $O_p$ of $p$ on which the indices $l$ with $p\in\partial T_l$ never affect the maximum, the indices with $p\notin\overline{T_l}$ do not occur, and the indices with $p\in T_l$ occur at every point. On $O_p$ the function $\psi$ is therefore the maximum of the $\widehat\varphi_l$ over the fixed finite set of indices $l$ with $p\in T_l$.

Let $W_0=\bigcup_{p\in Y}O_p$. A finite maximum of functions in $\PSH(O_p,(1-\varepsilon)\omega)$ belongs to this class, by the last remark following \cref{def:quasi-psh-and-theta-psh}. Hence $\psi\in\PSH(W_0,(1-\varepsilon)\omega)$; moreover $\psi$ is continuous on $W_0\setminus Y$, and $\mathrm e^\psi$, which is locally a finite maximum of the functions $\mathrm e^{\widehat\varphi_l}=\mathrm e^{\varphi_l}\mathrm e^{-\eta f_l}$, is H\"older continuous on $W_0$ by \cref{lem:manifold-holder-operations}(3). On $Y$ every $\widehat\varphi_l$ restricts to $\varphi-\eta f_l\leq\varphi$, while $\widehat\varphi_j=\varphi$ near $p$ for the index $j$ chosen above. Consequently $\psi|_Y=\varphi$.

\emph{Step 5. The global extension.} Choose an open neighborhood $G$ of $Y$ with $\overline G\subset W_0$, for instance a finite union of coordinate balls with closures in $W_0$ covering the compact set $Y$. Its boundary $\partial G$ is compact and disjoint from $Y$, and both $\psi$ and $F$ are finite and continuous near $\partial G$. Choose a constant $M$ with $M>\max_{\partial G}(\psi-F)$; since $F+M-\psi$ is continuous on $W_0\setminus Y$ and positive on $\partial G$, it is positive on an open neighborhood $G_1\subset W_0\setminus Y$ of $\partial G$. Set
\[
 \widetilde\varphi=
 \begin{cases}
  \max\{\psi,F+M\}&\text{on }G,\\
  F+M&\text{on }X\setminus G.
 \end{cases}
\]
The two expressions agree on $G_1$, so $\widetilde\varphi=F+M$ on the open set $G_1\cup(X\setminus\overline G)$, which contains $X\setminus G$. The function $F+M$ belongs to $\PSH(X,(1-4\varepsilon)\omega)\subset\PSH(X,(1-\varepsilon)\omega)$, and $\psi\in\PSH(W_0,(1-\varepsilon)\omega)$. Let $B$ be a coordinate ball and $h$ a smooth function on $B$ with $\ddc h=(1-\varepsilon)\omega$ (\cref{lem:smooth-potential-on-ball}). Then $F+M+h\in\PSH(B)$ and $\psi+h\in\PSH(G\cap B)$, and for every point $x$ of the boundary of $G\cap B$ in $B$, which lies in $\partial G$, we have $\limsup_{G\ni y\to x}(\psi+h)(y)\leq F(x)+M+h(x)$, because $\psi<F+M$ on $G_1$ and $F+M+h$ is continuous at $x$. By \cref{prop:psh-max-gluing}, $\widetilde\varphi+h$ is plurisubharmonic on $B$. Since such balls cover $X$, $\widetilde\varphi$ is quasi-plurisubharmonic and $(1-\varepsilon)\omega+\ddc\widetilde\varphi\geq0$, that is, $\widetilde\varphi\in\PSH(X,\omega)$ and $\omega+\ddc\widetilde\varphi\geq\varepsilon\omega$. Since $F=-\infty$ on $Y$, we have $\widetilde\varphi|_Y=\psi|_Y=\varphi$. The function $\widetilde\varphi$ is continuous on $X\setminus Y$, because $\psi$ and $F$ are continuous on $W_0\setminus Y$ and $X\setminus Y$, respectively. Its exponential equals $\max\{\mathrm e^\psi,\mathrm e^M\mathrm e^F\}$ on $G$ and $\mathrm e^M\mathrm e^F$ on the open set $G_1\cup(X\setminus\overline G)$; these two open sets cover $X$, so $\mathrm e^{\widetilde\varphi}$ is H\"older continuous on $X$ by \cref{lem:manifold-holder-operations}(3).

\emph{Step 6. Analytic singularities.} Finally suppose that $\varphi$ has analytic singularities (\cref{def:rational-analytic-singularities}). By \cref{prop:rational-analytic-operations}, so does the holomorphic pullback $\varphi\circ\pi_j$, which is not identically $-\infty$, and the function $F$ has analytic singularities by Step 1. By the same proposition, $\max\{\varphi\circ\pi_j,F\}$ has analytic singularities; adding the smooth function $AR$ changes the bounded remainder in its local presentations \eqref{eq:rational-analytic-singularities} by a locally bounded function, so $\max\{\varphi\circ\pi_j,F\}+AR$ has analytic singularities. By \cref{lem:background-logarithmic-sum},
\[
 \max\{\varphi\circ\pi_j,F\}+AR\leq\varphi_j\leq\max\{\varphi\circ\pi_j,F\}+AR+\log2.
\]
Hence a local presentation \eqref{eq:rational-analytic-singularities} of $\max\{\varphi\circ\pi_j,F\}+AR$ yields one of $\varphi_j$, with the same holomorphic functions and a bounded remainder changed by at most $\log2$; so $\varphi_j$ has analytic singularities. Subtracting the smooth function $\eta f_j$ changes the remainders by bounded functions, and the locally fixed finite maxima of Step 4 preserve analytic singularities by \cref{prop:rational-analytic-operations}; so $\psi$ has analytic singularities on $W_0$. On $G$, the function $\widetilde\varphi$ is the maximum of $\psi$ and $F+M$, and on the open set $G_1\cup(X\setminus\overline G)$ it equals $F+M$; both functions have analytic singularities. Therefore $\widetilde\varphi$ has analytic singularities on $X$.
\end{proof}
\index{K\"ahler current!submanifold potential extension}
\index{plurisubharmonic function!H\"older exponential extension}
